\scanpage{25}
\mspage{-2}
the map from $\handwritten{\otimes}^{p}\mathbb{A}$ to $\mathbb{A}$ which is defined by the multiplication $m$ of $A$.

\noindent - $\handwritten{\Sigma}^{p}\mathbb{A}$ is the subalgebra of $\handwritten{\otimes}^{p}\mathbb{A}$ which consists of the symmetric tensors of order $p$. Since the diagonal $\handwritten{\Delta}$ is commutative, $\handwritten{\Delta}^{p}$ is the composite of a map $a:\mathbb{A}\mathrel{\text{-{}-{}-{}-{}-}\!\handwritten{>}}\handwritten{\Sigma}^{p}\mathbb{A}$ and the canonical injection $i:\handwritten{\Sigma}^{p}\mathbb{A}\mathrel{\text{-{}-{}-{}-{}-}\!\handwritten{>}}\handwritten{\otimes}^{p}\mathbb{A}$.

If $(a_i)$ is a basis of $\mathbb{A}$, the products $b=a_{i_1}\handwritten{\otimes}\ldots\handwritten{\otimes}a_{i_r}$ form a basis of $\handwritten{\otimes}^{p}\mathbb{A}$. If $t_b$ is the subgroup of the symmetric group $s_p$ consisting of the permutations leaving $b$ fixed, we put $S(b)=\handwritten{\sum}\handwritten{\sigma}_{\deleted{\unsure{i}}j}(b)$, the $\handwritten{\sigma}_{\deleted{\unsure{i}}j}$ ranging over a system of representatives of $s_p\handwritten{/}t_b$. It is clear that the $S(b)$ form a basis of $\handwritten{\Sigma}^{p}\mathbb{A}$.

\noindent - $S^{p}\mathbb{A}$ is the image of $\handwritten{\otimes}^{p}\mathbb{A}$ in the symmetric algebra of $A$. Since the multiplicatipn $m$ is comm\deleted{j}utative, $m^p$ is the composite of the canonical surjection $q:\handwritten{\otimes}^{p}\mathbb{A}\mathrel{\text{-{}-{}-{}-{}-}\!\handwritten{>}}S^p\mathbb{A}$ and a ma\deleted{\unsure{s|x}}p $b:S^p\mathbb{A}\mathrel{\text{-{}-{}-{}-{}-{}-}\!\handwritten{>}}\mathbb{A}$.

If $(a_i)$ is a basis of $\mathbb{A}$, $S^p\mathbb{A}$ is the set of homogeneous polynomials of degree $p$ in the $a_i$.

\noindent - We are going to define a map $r:\handwritten{\Sigma}^{p}\mathbb{A}\mathrel{\text{-{}-{}-{}-{}-}\!\handwritten{>}}\mathbb{A}^{(p)}$ by exhibiting the $r(S(b))$ :

if $a_{i_1}=\ldots=a_{i_r}=a$, we put $r(S(b))=r(a\handwritten{\otimes}\ldots\handwritten{\otimes}a)=1\handwritten{\otimes}a$.

if the $a_{i_j}$ are not all equal, we put $r(S(b))=r(\handwritten{\sum}\handwritten{\sigma}_j(b))=0$.

\noindent - We likewise define a map $j:\mathbb{A}^{(p)}\mathrel{\text{-{}-{}-{}-{}-}\!\handwritten{>}}S^p\mathbb{A}$ such that
\[
j(1\handwritten{\otimes}a_i)=a_i^p=q(a\handwritten{\otimes}\ldots\handwritten{\otimes}a)\ .
\]
% archive sheet number: 24
